% ============================================================
%  يُترجَم بـ XeLaTeX  (في Overleaf: Menu → Compiler → XeLaTeX)
%  الخط: Amiri (متوفر في Overleaf، أو حمّله من github.com/aliftype/amiri)
% ============================================================
\documentclass[12pt,a4paper]{article}
\usepackage[a4paper,margin=1.7cm,top=2.3cm,bottom=2.9cm,headheight=15pt,footskip=1.5cm]{geometry}
\usepackage{amsmath,amssymb}
\usepackage{icomma}
\usepackage[english,bidi=default,layout=lists]{babel}
\babelprovide[import,main]{arabic}
\babelfont{rm}[Scale=1.08]{Amiri}
\babelfont{sf}[Scale=1.08]{Amiri}
\usepackage[most]{tcolorbox}
\usepackage{tikz}
\usetikzlibrary{arrows.meta,calc,shadings}
\usepackage{array}
\usepackage{colortbl}
\usepackage{fancyhdr}
\usepackage{titlesec}
\usepackage{enumitem}
\usepackage{collcell}
\usepackage{needspace}
\usepackage{etoolbox}
\newcolumntype{C}{>{\collectcell\ar}c<{\endcollectcell}}
\pretocmd{\section}{\Needspace{15\baselineskip}}{}{}
\pretocmd{\subsection}{\Needspace{11\baselineskip}}{}{}

\definecolor{cDef}{HTML}{1F4E9C}
\definecolor{cThm}{HTML}{B3261E}
\definecolor{cRem}{HTML}{2E7D32}
\definecolor{cEx}{HTML}{C76B00}
\definecolor{cMeth}{HTML}{6A1B9A}
\definecolor{cSol}{HTML}{00695C}
\definecolor{cSec}{HTML}{1F3A68}

\titleformat{\section}{\Large\bfseries\color{cSec}}{\thesection.}{0.6em}{}[{\color{cSec!60}\titlerule[1pt]}]
\titleformat{\subsection}{\large\bfseries\color{cThm!85!black}}{\thesubsection.}{0.5em}{}
\titlespacing*{\section}{0pt}{14pt}{8pt}
\titlespacing*{\subsection}{0pt}{10pt}{5pt}
\setlength{\parindent}{0pt}
\setlength{\parskip}{4pt}
\setlist{itemsep=2pt,topsep=3pt}

\newtcolorbox{arbox}[3][]{enhanced,
  colback=#2!5,colframe=#2!80!black,arc=2mm,boxrule=0.8pt,
  left=6pt,right=6pt,top=8pt,bottom=5pt,
  fonttitle=\bfseries,coltitle=white,
  title={\foreignlanguage{arabic}{#3}},
  attach boxed title to top right={xshift=-5mm,yshift=-3mm},
  boxed title style={colback=#2!80!black,arc=1.5mm,boxrule=0pt},
  before upper={\selectlanguage{arabic}},
  before skip=8pt,after skip=8pt,breakable,#1}
\newenvironment{BOX}[2]{\begin{otherlanguage}{english}\begin{arbox}{#1}{#2}}{\end{arbox}\end{otherlanguage}}
\newenvironment{tarif}[1][تعريف]{\begin{BOX}{cDef}{#1}}{\end{BOX}}
\newenvironment{khassia}[1][خاصية]{\begin{BOX}{cThm}{#1}}{\end{BOX}}
\newenvironment{molahada}[1][ملاحظة]{\begin{BOX}{cRem}{#1}}{\end{BOX}}
\newenvironment{mithal}[1][مثال]{\begin{BOX}{cEx}{#1}}{\end{BOX}}
\newenvironment{tariqa}[1][طريقة]{\begin{BOX}{cMeth}{#1}}{\end{BOX}}

\newcounter{exo}
\newenvironment{tamrin}[1]{\stepcounter{exo}\begin{BOX}{cSec}{التمرين \theexo: #1}}{\end{BOX}}
\newenvironment{hal}{\Needspace{12\baselineskip}\begin{otherlanguage}{english}\begin{arbox}[breakable]{cSol}{حل التمرين \theexo}}{\end{arbox}\end{otherlanguage}}

\newenvironment{LTR}{\begin{otherlanguage}{english}}{\end{otherlanguage}}
\newcommand{\ar}[1]{\foreignlanguage{arabic}{#1}}

\renewcommand{\dots}{\mathinner{\ldotp\ldotp\ldotp}}
\renewcommand{\ldots}{\mathinner{\ldotp\ldotp\ldotp}}
\renewcommand{\cdots}{\mathinner{\cdotp\cdotp\cdotp}}
\newcommand{\N}{\mathbb{N}}
\newcommand{\Z}{\mathbb{Z}}
\newcommand{\D}{\mathbb{D}}
\newcommand{\Q}{\mathbb{Q}}
\newcommand{\R}{\mathbb{R}}
\newcommand{\pgcd}{\operatorname{PGCD}}
\newcommand{\ppcm}{\operatorname{PPCM}}

% \NLseg{a}{b}{[}{]} ; \NLright{a}{[} ; \NLleft{b}{]}  (interval on a number line)
\newcommand{\NLseg}[4]{\begin{LTR}\begin{tikzpicture}[baseline=-0.6ex]
 \draw[-{Stealth[length=2mm]},gray!70](0,0)--(3.4,0);
 \draw[line width=2pt,cThm!80!black](0.7,0)--(2.7,0);
 \node[cThm!80!black,font=\LARGE] at (0.7,0){$#3$};
 \node[cThm!80!black,font=\LARGE] at (2.7,0){$#4$};
 \node[below=3pt,font=\footnotesize] at (0.7,0){$#1$};
 \node[below=3pt,font=\footnotesize] at (2.7,0){$#2$};
\end{tikzpicture}\end{LTR}}
\newcommand{\NLright}[2]{\begin{LTR}\begin{tikzpicture}[baseline=-0.6ex]
 \draw[-{Stealth[length=2mm]},gray!70](0,0)--(3.4,0);
 \draw[line width=2pt,cThm!80!black](1.3,0)--(3.3,0);
 \node[cThm!80!black,font=\LARGE] at (1.3,0){$#2$};
 \node[below=3pt,font=\footnotesize] at (1.3,0){$#1$};
\end{tikzpicture}\end{LTR}}
\newcommand{\NLleft}[2]{\begin{LTR}\begin{tikzpicture}[baseline=-0.6ex]
 \draw[-{Stealth[length=2mm]},gray!70](0,0)--(3.4,0);
 \draw[line width=2pt,cThm!80!black](0.05,0)--(2.1,0);
 \node[cThm!80!black,font=\LARGE] at (2.1,0){$#2$};
 \node[below=3pt,font=\footnotesize] at (2.1,0){$#1$};
\end{tikzpicture}\end{LTR}}

% ---------- social-media icons (pure TikZ) ----------
\definecolor{cFB}{HTML}{1877F2}
\definecolor{cTG}{HTML}{229ED9}
\definecolor{cIG1}{HTML}{FEDA75}
\definecolor{cIG2}{HTML}{D62976}
\definecolor{cIG3}{HTML}{4F5BD5}
\newcommand{\fbicon}{\begin{LTR}\begin{tikzpicture}[baseline=-0.75ex,x=1em,y=1em]
 \fill[cFB] (0,0) circle (0.55);
 \node[white,inner sep=0,font=\usefont{TU}{lmss}{bx}{n}\small] at (0.04,-0.02) {f};
\end{tikzpicture}\end{LTR}}
\newcommand{\igicon}{\begin{LTR}\begin{tikzpicture}[baseline=-0.75ex,x=1em,y=1em]
 \shade[lower left=cIG1,lower right=cIG2,upper left=cIG2,upper right=cIG3,rounded corners=0.28em] (-0.55,-0.55) rectangle (0.55,0.55);
 \draw[white,line width=0.09em,rounded corners=0.18em] (-0.36,-0.36) rectangle (0.36,0.36);
 \draw[white,line width=0.09em] (0,0) circle (0.17);
 \fill[white] (0.22,0.22) circle (0.05);
\end{tikzpicture}\end{LTR}}
\newcommand{\tgicon}{\begin{LTR}\begin{tikzpicture}[baseline=-0.75ex,x=1em,y=1em]
 \fill[cTG] (0,0) circle (0.55);
 \fill[white] (-0.33,0.0) -- (0.32,0.27) -- (0.2,-0.3) -- (0.0,-0.15) -- (-0.08,-0.28) -- (-0.08,-0.08) -- cycle;
 \fill[cTG!70!black] (-0.08,-0.08) -- (0.25,0.2) -- (0.0,-0.15) -- (-0.08,-0.28) -- cycle;
\end{tikzpicture}\end{LTR}}
% page number: black box, white two-digit number, at the right of the social row
\newcommand{\PageBox}{\begin{LTR}\tikz[baseline=(p.base)]\node[fill=black,text=white,inner sep=2.5pt,minimum width=1.6em,minimum height=1.3em,font=\usefont{TU}{lmss}{bx}{n}\small](p){\ifnum\value{page}<10 0\fi\arabic{page}};\end{LTR}}
% visual left-to-right order: Telegram | Instagram | Facebook | page box
\newcommand{\SocialRow}{\begin{otherlanguage}{english}%
 \ar{\textbf{الأستاذ سوايسية هشام للرياضيات}}\ \tgicon\hfill
 \textbf{prof\_hicham\_math}\ \igicon\hfill
 \ar{\textbf{الأستاذ سوايسية هشام للرياضيات}}\ \fbicon\hspace{2.2em}\PageBox\end{otherlanguage}}

\pagestyle{fancy}
\fancyhf{}
\renewcommand{\headrulewidth}{0.6pt}
\renewcommand{\footrulewidth}{0.4pt}
% \SetHeader{right head: document title}{left head: level -- stream}
\newcommand{\SetHeader}[2]{%
 \fancyhead[C]{\parbox{\textwidth}{\ar{#1\hfill #2}}}%
 \fancyfoot[C]{\parbox{\textwidth}{\small\SocialRow}}}
\usepackage{tkz-tab}
\usepackage{multicol}
% variation / sign table (LTR)
\newenvironment{TAB}{\par\begin{center}\begin{LTR}\begin{tikzpicture}[scale=0.95,every node/.style={font=\small}]}{\end{tikzpicture}\end{LTR}\end{center}}
% grid + axes, clipped: \begin{GR}{xmin}{xmax}{ymin}{ymax}{unit cm}
\newenvironment{GR}[5]{\par\begin{center}\begin{LTR}\begin{tikzpicture}[x=#5cm,y=#5cm,font=\footnotesize]
 \draw[gray!25,very thin] (#1,#3) grid (#2,#4);
 \draw[-{Stealth[length=2mm]}] (#1,0)--(#2+0.3,0) node[below left]{$x$};
 \draw[-{Stealth[length=2mm]}] (0,#3)--(0,#4+0.3) node[below left]{$y$};
 \node[below left] at (0,0){$O$};
 \begin{scope}\clip (#1,#3) rectangle (#2,#4);}
{\end{scope}\end{tikzpicture}\end{LTR}\end{center}}
\newcommand{\Ticks}[2]{\foreach \t in {#1} \node[below,font=\scriptsize] at (\t,0){$\t$}; \foreach \t in {#2} \node[left,font=\scriptsize] at (0,\t){$\t$};}
\newcommand{\bx}[1]{\ensuremath{\colorbox{cDef!12}{$\displaystyle #1$}}}
\newcommand{\TwoCol}[4]{\par\begin{LTR}\noindent\begin{minipage}[t]{#1\textwidth}\selectlanguage{arabic}#2\end{minipage}\hfill\begin{minipage}[t]{#3\textwidth}\selectlanguage{arabic}#4\end{minipage}\end{LTR}\par}
\newcommand{\RL}[4]{\TwoCol{#3}{#4}{#1}{#2}}  % \RL{w1}{right content}{w2}{left content}
\newcommand{\ok}{$\textcolor{cRem}{\checkmark}$}
\newcommand{\limn}{\displaystyle\lim_{n\to+\infty}}
% sub-questions, each on its own line
\newenvironment{sq}{\begin{itemize}[leftmargin=2.4em,itemsep=3pt]}{\end{itemize}}
\newcommand{\qa}{\item[(أ)]}\newcommand{\qb}{\item[(ب)]}\newcommand{\qc}{\item[(جـ)]}\newcommand{\qd}{\item[(د)]}\newcommand{\qe}{\item[(هـ)]}
\newcommand{\pa}{\par\smallskip\textbf{(أ)} }\newcommand{\pb}{\par\smallskip\textbf{(ب)} }\newcommand{\pc}{\par\smallskip\textbf{(جـ)} }\newcommand{\pd}{\par\smallskip\textbf{(د)} }\newcommand{\pe}{\par\smallskip\textbf{(هـ)} }
\newcommand{\QQ}[1]{\par\smallskip\textbf{#1.} }
\newenvironment{tmr}{\stepcounter{exo}\begin{BOX}{cSec}{تمرين محلول \theexo}}{\end{BOX}}
\newcommand{\ABC}[3]{\\[2pt]\begin{tabular}{p{0.28\textwidth}p{0.28\textwidth}p{0.28\textwidth}}(جـ) \ #3 & (ب) \ #2 & (أ) \ #1\end{tabular}}
\pretocmd{\section}{\Needspace{24\baselineskip}}{}{}
\newtcbox{\bxtbox}{on line,colback=cDef!12,colframe=cDef!12,boxrule=0pt,arc=1pt,boxsep=0pt,left=2pt,right=2pt,top=1pt,bottom=1pt}
\newcommand{\bxt}[1]{\foreignlanguage{english}{\bxtbox{\ar{#1}}}}
\usepackage{pgfplots}
\pgfplotsset{compat=1.16}
\pgfplotsset{sc/.style={width=0.8\textwidth,height=6.3cm,axis lines=left,grid=both,minor tick num=1,
  major grid style={gray!40},minor grid style={gray!15},tick label style={font=\scriptsize},label style={font=\footnotesize},
  /pgf/number format/use comma,/pgf/number format/1000 sep={},scaled ticks=false,
  xlabel style={at={(ticklabel* cs:1)},anchor=north east},ylabel style={rotate=-90,at={(ticklabel* cs:1)},anchor=south west}},
  pt/.style={only marks,mark=*,mark size=2.2pt,color=cDef},
  gp/.style={only marks,mark=square*,mark size=2.6pt,color=cThm},
  gq/.style={only marks,mark=triangle*,mark size=3pt,color=cEx},
  ln/.style={thick,color=cThm,samples=2},
  lm/.style={thick,dashed,color=cRem,samples=2}}
% scatter plot on graph paper: \begin{SC}{axis options}
\newenvironment{SC}[1]{\par\begin{center}\begin{LTR}\begin{tikzpicture}\begin{axis}[sc,#1]}{\end{axis}\end{tikzpicture}\end{LTR}\end{center}}
\SetHeader{\small\bfseries\color{cSec}تمارين تجريبية: الإحصاء}{\small\color{cSec}السنة الثالثة ثانوي -- تسيير واقتصاد}
\BeforeBeginEnvironment{tamrin}{\Needspace{12\baselineskip}}
\BeforeBeginEnvironment{hal}{\Needspace{20\baselineskip}}
\begin{document}
\begin{LTR}
\begin{tcolorbox}[enhanced,colback=cSec,colframe=cSec,arc=3mm,halign=center,boxrule=0pt,top=12pt,bottom=12pt,before upper={\selectlanguage{arabic}}]
{\color{white}\fontsize{24}{32}\selectfont\bfseries تمارين تجريبية: الإحصاء}\\[8pt]
{\color{white}\large سحابة النقط -- طريقة ماير -- تعديل بياني ومستقيم معطى -- تغيير المتغير}\\[5pt]
{\color{white}\large 11 تمرينا متدرجا مع الحلول المفصلة خطوة بخطوة}\\[5pt]{\color{white}\large السنة الثالثة ثانوي -- تسيير واقتصاد}
\end{tcolorbox}
\end{LTR}

\begin{tarif}[تذكير: القوانين التي نستعملها]
$(x_i;y_i)$ سلسلة إحصائية ذات متغيرين عدد نقطها $n$.
\begin{itemize}
\item \textbf{النقطة المتوسطة:} $G(\bar x;\bar y)$ حيث $\bar x=\dfrac{1}{n}\sum x_i$ و $\bar y=\dfrac{1}{n}\sum y_i$.
\item \textbf{تعديل بياني:} نرسم مستقيما يشمل $G$ ويترك تقريبا نفس عدد النقط من كل جهة، ثم نعيّن معادلته من $G$ ونقطة ثانية منه.
\item \textbf{مستقيم تعديل معطى:} نتحقق أنه يشمل $G$ أي $\bar y=a\bar x+b$ (بتقريب التدوير)، ثم نرسمه بنقطتين.
\item \textbf{طريقة ماير:} نرتب النقط حسب $x_i$، نقسمها إلى مجموعتين متساويتين في العدد، نحسب النقطتين المتوسطتين $G_1$ و $G_2$؛ المستقيم $(G_1G_2)$ هو مستقيم ماير ويشمل $G$.
\item \textbf{تغيير المتغير:} إذا كان $z=\ln y$ و $z=ax+b$ فإن $y=e^{ax+b}=e^{b}\left(e^{a}\right)^{x}$، أي $y=k\,q^{x}$ مع $k=e^{b}$ و $q=e^{a}$.
\end{itemize}
\end{tarif}

\section{سحابة النقط، النقطة المتوسطة وطريقة ماير}
% ------------------------------------------------------------------
\begin{tamrin}{مبيعات متجر -- تعديل بطريقة ماير}
يعطي الجدول التالي المبيعات الشهرية $y_i$ (بالألف دينار جزائري) لمتجر خلال الأشهر الستة الأولى من سنة 2025، حيث $x_i$ رتبة الشهر ($x_i=1$ يوافق جانفي).
\begin{center}\begin{LTR}\renewcommand{\arraystretch}{1.35}\small
\begin{tabular}{|c|c|c|c|c|c|c|}\hline
\rowcolor{cSec!12}\ar{رتبة الشهر} $x_i$ & $1$ & $2$ & $3$ & $4$ & $5$ & $6$\\\hline
\ar{المبيعات} $y_i$ & $420$ & $470$ & $490$ & $520$ & $590$ & $630$\\\hline
\end{tabular}\end{LTR}\end{center}
\QQ{1} مثّل سحابة النقط $M_i(x_i;y_i)$ في معلم متعامد (على محور الفواصل: $1\,\text{cm}$ لكل شهر؛ وعلى محور التراتيب: $1\,\text{cm}$ لكل $20$ ألف دينار، ونبدأ التدريج من $400$).
\QQ{2} احسب إحداثيي النقطة المتوسطة $G$ لهذه السلسلة ثم علّمها.
\QQ{3} هل يبرّر شكل السحابة إجراء تعديل خطي؟
\QQ{4} نقسم السلسلة إلى مجموعتين: الأولى تضم النقط الثلاث الأولى، والثانية تضم النقط الثلاث الأخيرة.
\begin{sq}
\qa احسب إحداثيي $G_1$ و $G_2$ النقطتين المتوسطتين للمجموعتين.
\qb عيّن معادلة مستقيم ماير $(G_1G_2)$ وتحقق أن $G$ تنتمي إليه، ثم ارسمه.
\end{sq}
\QQ{5} باستعمال هذا التعديل:
\begin{sq}
\qa قدّر مبيعات شهر سبتمبر 2025.
\qb ابتداء من أي شهر من سنة 2025 تتجاوز المبيعات $800$ ألف دينار (إذا استمر الاتجاه نفسه)؟
\end{sq}
\end{tamrin}
\begin{hal}
\QQ{1}{\bfseries و 2 و 4.} سحابة النقط (بالأزرق)، النقطة المتوسطة $G$ (بالأحمر)، النقطتان $G_1$ و $G_2$ (بالبرتقالي) ومستقيم ماير:
\begin{SC}{xmin=0,xmax=7.4,ymin=400,ymax=660,xtick={0,...,7},ytick={400,440,...,640},xlabel={$x_i$},ylabel={$y_i$}}
\addplot[ln,domain=0:7.2]{40*x+380};
\addplot[pt] coordinates {(1,420)(2,470)(3,490)(4,520)(5,590)(6,630)};
\addplot[gq] coordinates {(2,460)(5,580)};
\addplot[gp] coordinates {(3.5,520)};
\node[above left] at (axis cs:3.5,520){$G$};
\node[below right] at (axis cs:2,460){$G_1$};
\node[below right] at (axis cs:5,580){$G_2$};
\node[cThm,above left] at (axis cs:6.9,656){$(G_1G_2)$};
\end{SC}
\QQ{2} عدد النقط $n=6$:
\[\bar x=\frac{1+2+3+4+5+6}{6}=\frac{21}{6}=3{,}5\]\[\bar y=\frac{420+470+490+520+590+630}{6}=\frac{3120}{6}=520\]
إذن \bx{G(3{,}5\,;\,520)}

\QQ{3} نقط السحابة متقاربة من مستقيم (تتزايد المبيعات بانتظام تقريبا)، إذن يمكن إجراء \textbf{تعديل خطي}.

\QQ{4}\pa المجموعة الأولى: النقط ذات الفواصل $1$ ، $2$ ، $3$:
\[x_{G_1}=\frac{1+2+3}{3}=2\qquad y_{G_1}=\frac{420+470+490}{3}=\frac{1380}{3}=460\qquad\text{إذن}\qquad \bx{G_1(2\,;\,460)}\]
المجموعة الثانية: النقط ذات الفواصل $4$ ، $5$ ، $6$:
\[x_{G_2}=\frac{4+5+6}{3}=5\qquad y_{G_2}=\frac{520+590+630}{3}=\frac{1740}{3}=580\qquad\text{إذن}\qquad \bx{G_2(5\,;\,580)}\]
\pb معادلة $(G_1G_2)$ من الشكل $y=ax+b$ حيث المعامل الموجه:
\[a=\frac{y_{G_2}-y_{G_1}}{x_{G_2}-x_{G_1}}=\frac{580-460}{5-2}=\frac{120}{3}=40\]
و $G_1\in(G_1G_2)$ ومنه $460=40\times2+b$ أي $b=380$. إذن معادلة مستقيم ماير: \bx{y=40x+380}

التحقق: $40\times3{,}5+380=140+380=520=\bar y$، إذن \bx{G\in(G_1G_2)} \ok

\textbf{تفسير:} $a=40$ يعني أن المبيعات تزداد في المتوسط بـ $40$ ألف دينار كل شهر.

\QQ{5}\pa شهر سبتمبر رتبته $x=9$: \ $y=40\times9+380=740$. تقدَّر مبيعات سبتمبر بـ \bxt{$740$ ألف دينار} تقريبا.
\pb نحل المتراجحة: $40x+380>800$ \ أي \ $40x>420$ \ أي \ $x>10{,}5$. أصغر رتبة طبيعية هي $x=11$، إذن ابتداء من \bxt{شهر نوفمبر 2025} تتجاوز المبيعات $800$ ألف دينار.
\end{hal}

% ------------------------------------------------------------------
\begin{tamrin}{السعر والطلب -- مستقيم ماير والإيراد الأعظمي}
درست مؤسسة الطلب اليومي $y_i$ (عدد الوحدات المطلوبة) على منتوج بدلالة سعر الوحدة $x_i$ (بالدينار) فتحصلت على:
\begin{center}\begin{LTR}\renewcommand{\arraystretch}{1.35}\small
\begin{tabular}{|c|c|c|c|c|c|c|}\hline
\rowcolor{cSec!12}\ar{السعر} $x_i$ & $20$ & $30$ & $40$ & $50$ & $60$ & $70$\\\hline
\ar{الطلب} $y_i$ & $412$ & $381$ & $350$ & $302$ & $275$ & $251$\\\hline
\end{tabular}\end{LTR}\end{center}
\QQ{1} مثّل سحابة النقط $M_i(x_i;y_i)$ (على محور الفواصل: $1\,\text{cm}$ لكل $10$ دنانير؛ وعلى محور التراتيب: $1\,\text{cm}$ لكل $20$ وحدة ابتداء من $200$).
\QQ{2} احسب إحداثيي النقطة المتوسطة $G$.
\QQ{3} باستعمال طريقة ماير (مجموعة أولى: النقط الثلاث الأولى) بيّن أن معادلة مستقيم التعديل هي: $y=-3{,}5x+486$.
\QQ{4} فسّر اقتصاديا المعامل $-3{,}5$.
\QQ{5}
\begin{sq}
\qa قدّر الطلب اليومي إذا كان سعر الوحدة $80$ دينارا.
\qb ما السعر الذي يوافق طلبا يوميا قدره $300$ وحدة؟ (دوّر إلى $10^{-2}$).
\end{sq}
\QQ{6} نعتبر أن الطلب معطى بـ $y=-3{,}5x+486$ من أجل $x\in[20;100]$. الإيراد اليومي هو $R(x)=x\cdot y$.
\begin{sq}
\qa بيّن أن $R(x)=-3{,}5x^{2}+486x$.
\qb ادرس اتجاه تغير الدالة $R$ على $[20;100]$، واستنتج السعر الذي يجعل الإيراد أعظميا وقيمة هذا الإيراد (دوّر إلى $10^{-2}$).
\end{sq}
\end{tamrin}
\begin{hal}
\QQ{1} السحابة ومستقيم ماير ($G$ بالأحمر و $G_1,G_2$ بالبرتقالي):
\begin{SC}{xmin=0,xmax=88,ymin=200,ymax=440,xtick={0,10,...,80},ytick={200,240,...,440},xlabel={$x_i$},ylabel={$y_i$}}
\addplot[ln,domain=0:85]{-3.5*x+486};
\addplot[pt] coordinates {(20,412)(30,381)(40,350)(50,302)(60,275)(70,251)};
\addplot[gq] coordinates {(30,381)(60,276)};
\addplot[gp] coordinates {(45,328.5)};
\node[above right] at (axis cs:45,328.5){$G$};
\node[below left] at (axis cs:30,381){$G_1$};
\node[below left] at (axis cs:60,276){$G_2$};
\end{SC}
\QQ{2}
\[\bar x=\frac{20+30+40+50+60+70}{6}=\frac{270}{6}=45\]\[\bar y=\frac{412+381+350+302+275+251}{6}=\frac{1971}{6}=328{,}5\]
إذن \bx{G(45\,;\,328{,}5)}

\QQ{3} المجموعة الأولى ($x_i=20;30;40$) والثانية ($x_i=50;60;70$):
\[G_1\left(\frac{20+30+40}{3}\,;\,\frac{412+381+350}{3}\right)=G_1\left(30\,;\,\frac{1143}{3}\right)=G_1(30\,;\,381)\]
\[G_2\left(\frac{50+60+70}{3}\,;\,\frac{302+275+251}{3}\right)=G_2\left(60\,;\,\frac{828}{3}\right)=G_2(60\,;\,276)\]
\[a=\frac{276-381}{60-30}=\frac{-105}{30}=-3{,}5\qquad b=y_{G_1}-a\,x_{G_1}=381-(-3{,}5)\times30=381+105=486\]
إذن \bx{(G_1G_2):\ y=-3{,}5x+486}. \ التحقق: $-3{,}5\times45+486=-157{,}5+486=328{,}5=\bar y$ ومنه $G\in(G_1G_2)$ \ok

\QQ{4} كلما ارتفع سعر الوحدة بدينار واحد \textbf{انخفض} الطلب اليومي بـ $3{,}5$ وحدات تقريبا (المعامل سالب: الطلب دالة متناقصة للسعر).

\QQ{5}\pa $y=-3{,}5\times80+486=-280+486=206$. الطلب المقدَّر: \bxt{$206$ وحدة يوميا}
\pb $-3{,}5x+486=300$ \ أي \ $3{,}5x=186$ \ أي \ $x=\dfrac{186}{3{,}5}\approx53{,}142\ldots$ \ إذن السعر \bx{x\approx53{,}14\ \text{DA}}

\QQ{6}\pa $R(x)=x\left(-3{,}5x+486\right)=\bx{-3{,}5x^{2}+486x}$
\pb $R'(x)=-7x+486$، و $R'(x)=0$ يكافئ $x=\dfrac{486}{7}\approx69{,}43$.
لما $x<\dfrac{486}{7}$ يكون $R'(x)>0$، ولما $x>\dfrac{486}{7}$ يكون $R'(x)<0$:
\begin{TAB}
\tkzTabInit[lgt=1.6,espcl=3.2]{$x$/1,$R'(x)$/1,$R(x)$/2}{$20$,$\dfrac{486}{7}$,$100$}
\tkzTabLine{,+,z,-,}
\tkzTabVar{-/$8320$,+/$\approx16871{,}14$,-/$13600$}
\end{TAB}
$R(20)=-1400+9720=8320$ ، \ $R(100)=-35000+48600=13600$ ، و:
\[R\left(\frac{486}{7}\right)=-3{,}5\times\frac{486^{2}}{49}+\frac{486^{2}}{7}=\frac{486^{2}}{14}=\frac{236196}{14}\approx16871{,}14\]

يكون الإيراد أعظميا لما \bx{x\approx69{,}43\ \text{DA}}، والإيراد الأعظمي \bx{R_{\max}\approx16871{,}14\ \text{DA}} يوميا.
\end{hal}

% ------------------------------------------------------------------
\begin{tamrin}{تطور عدد سكان مدينة -- ماير بثماني نقط}
يمثل الجدول التالي عدد سكان مدينة $y_i$ (بالآلاف) بين سنتي 2016 و 2023، حيث $x_i$ رتبة السنة.
\begin{center}\begin{LTR}\renewcommand{\arraystretch}{1.35}\small
\begin{tabular}{|c|c|c|c|c|c|c|c|c|}\hline
\rowcolor{cSec!12}\ar{السنة} & $2016$ & $2017$ & $2018$ & $2019$ & $2020$ & $2021$ & $2022$ & $2023$\\\hline
\ar{الرتبة} $x_i$ & $1$ & $2$ & $3$ & $4$ & $5$ & $6$ & $7$ & $8$\\\hline
\ar{السكان} $y_i$ & $24,2$ & $24,9$ & $25,7$ & $26,4$ & $27,5$ & $28,6$ & $29,4$ & $30,1$\\\hline
\end{tabular}\end{LTR}\end{center}
\QQ{1} مثّل سحابة النقط $M_i(x_i;y_i)$ (على محور الفواصل: $1\,\text{cm}$ لكل سنة؛ وعلى محور التراتيب: $1\,\text{cm}$ لكل ألف نسمة ابتداء من $24$).
\QQ{2} احسب إحداثيي النقطة المتوسطة $G$.
\QQ{3} عيّن معادلة مستقيم ماير $(G_1G_2)$، حيث $G_1$ و $G_2$ النقطتان المتوسطتان للنقط الأربع الأولى وللنقط الأربع الأخيرة على الترتيب، ثم تحقق أن $G\in(G_1G_2)$.
\QQ{4} فسّر المعامل الموجه لهذا المستقيم.
\QQ{5}
\begin{sq}
\qa قدّر عدد سكان المدينة سنة 2026.
\qb في أي سنة يتجاوز عدد السكان $35$ ألف نسمة لأول مرة؟
\end{sq}
\QQ{6} احسب النسبة المئوية لزيادة عدد السكان بين سنتي 2016 و 2023 (دوّر إلى $10^{-2}$).
\end{tamrin}
\begin{hal}
\QQ{1} السحابة ومستقيم ماير:
\begin{SC}{xmin=0,xmax=9.3,ymin=24,ymax=31,xtick={0,...,9},ytick={24,...,31},xlabel={$x_i$},ylabel={$y_i$}}
\addplot[ln,domain=0:9.2]{0.9*x+23.05};
\addplot[pt] coordinates {(1,24.2)(2,24.9)(3,25.7)(4,26.4)(5,27.5)(6,28.6)(7,29.4)(8,30.1)};
\addplot[gq] coordinates {(2.5,25.3)(6.5,28.9)};
\addplot[gp] coordinates {(4.5,27.1)};
\node[above left] at (axis cs:4.5,27.1){$G$};
\node[below right] at (axis cs:2.5,25.3){$G_1$};
\node[below right] at (axis cs:6.5,28.9){$G_2$};
\end{SC}
\QQ{2} $n=8$:
\[\bar x=\frac{1+2+\cdots+8}{8}=\frac{36}{8}=4{,}5\]
\[\bar y=\frac{24{,}2+24{,}9+25{,}7+26{,}4+27{,}5+28{,}6+29{,}4+30{,}1}{8}=\frac{216{,}8}{8}=27{,}1\]
إذن \bx{G(4{,}5\,;\,27{,}1)}

\QQ{3} المجموعة الأولى ($x_i=1;2;3;4$):
\[x_{G_1}=\frac{1+2+3+4}{4}=2{,}5\qquad y_{G_1}=\frac{24{,}2+24{,}9+25{,}7+26{,}4}{4}=\frac{101{,}2}{4}=25{,}3\]
إذن \bx{G_1(2{,}5\,;\,25{,}3)}. \ المجموعة الثانية ($x_i=5;6;7;8$):
\[x_{G_2}=\frac{5+6+7+8}{4}=6{,}5\qquad y_{G_2}=\frac{27{,}5+28{,}6+29{,}4+30{,}1}{4}=\frac{115{,}6}{4}=28{,}9\]
إذن \bx{G_2(6{,}5\,;\,28{,}9)}
\[a=\frac{28{,}9-25{,}3}{6{,}5-2{,}5}=\frac{3{,}6}{4}=0{,}9\qquad b=25{,}3-0{,}9\times2{,}5=25{,}3-2{,}25=23{,}05\]
إذن \bx{(G_1G_2):\ y=0{,}9x+23{,}05}. \ التحقق: $0{,}9\times4{,}5+23{,}05=4{,}05+23{,}05=27{,}1=\bar y$ \ok

\QQ{4} $a=0{,}9$: يزداد عدد السكان في المتوسط بـ $0{,}9$ ألف، أي بـ $900$ نسمة كل سنة.

\QQ{5}\pa سنة 2026 رتبتها $x=2026-2015=11$: \ $y=0{,}9\times11+23{,}05=32{,}95$. أي حوالي \bxt{$32\,950$ نسمة}
\pb $0{,}9x+23{,}05>35$ \ أي \ $x>\dfrac{11{,}95}{0{,}9}\approx13{,}28$، ومنه $x=14$، أي السنة $2015+14$: \bxt{سنة 2029}

\QQ{6} النسبة المئوية للزيادة: \ $\dfrac{30{,}1-24{,}2}{24{,}2}\times100=\dfrac{5{,}9}{24{,}2}\times100\approx\bx{24{,}38\,\%}$
\end{hal}

\section{مستقيم تعديل معطى، تعديل بياني، وماير في سياقات اقتصادية}
% ------------------------------------------------------------------
\begin{tamrin}{تكلفة الإنتاج -- مستقيم تعديل معطى}
يعطي الجدول التالي التكلفة الإجمالية $y_i$ (بالألف دينار) لإنتاج $x_i$ مئة وحدة من سلعة في ورشة:
\begin{center}\begin{LTR}\renewcommand{\arraystretch}{1.35}\small
\begin{tabular}{|c|c|c|c|c|c|}\hline
\rowcolor{cSec!12}\ar{الإنتاج} $x_i$ & $1$ & $2$ & $3$ & $4$ & $5$\\\hline
\ar{التكلفة} $y_i$ & $120$ & $150$ & $190$ & $230$ & $260$\\\hline
\end{tabular}\end{LTR}\end{center}
\QQ{1} مثّل سحابة النقط $M_i(x_i;y_i)$ (على محور الفواصل: $2\,\text{cm}$ لكل مئة وحدة؛ وعلى محور التراتيب: $1\,\text{cm}$ لكل $20$ ألف دينار ابتداء من $80$)، ثم احسب إحداثيي النقطة المتوسطة $G$.
\QQ{2} هل يبرّر شكل السحابة تعديلا خطيا؟
\QQ{3} نقبل أن المستقيم $(D):\ y=36x+82$ مستقيم تعديل لهذه السحابة. تحقق أن $G\in(D)$ ثم ارسم $(D)$.
\QQ{4} فسّر اقتصاديا العددين $36$ و $82$.
\QQ{5}
\begin{sq}
\qa قدّر تكلفة إنتاج $700$ وحدة.
\qb ميزانية الورشة $400$ ألف دينار. ما أكبر عدد من الوحدات يمكن إنتاجه؟
\end{sq}
\end{tamrin}
\begin{hal}
\QQ{1} السحابة ومستقيم التعديل $(D)$:
\begin{SC}{xmin=0,xmax=6.4,ymin=80,ymax=300,xtick={0,...,6},ytick={80,120,...,280},xlabel={$x_i$},ylabel={$y_i$}}
\addplot[ln,domain=0:6.2]{36*x+82};
\addplot[pt] coordinates {(1,120)(2,150)(3,190)(4,230)(5,260)};
\addplot[gp] coordinates {(3,190)};
\node[above left] at (axis cs:3,190){$G$};
\node[cThm,below right] at (axis cs:5.2,269){$(D)$};
\end{SC}
$\bar x=\dfrac{15}{5}=3$ \ و \ $\bar y=\dfrac{120+150+190+230+260}{5}=\dfrac{950}{5}=190$ \ إذن \bx{G(3\,;\,190)}

\QQ{2} النقط متقاربة من مستقيم (التكلفة تزداد بانتظام تقريبا): التعديل الخطي \textbf{مبرَّر}.

\QQ{3} $36\times3+82=108+82=190=\bar y$، إذن \bx{G\in(D)} \ok. للرسم: $G(3;190)$ والنقطة $(0;82)$.

\QQ{4} $82$: \textbf{التكاليف الثابتة} (لما $x=0$) حوالي $82$ ألف دينار. \ $36$: كل مئة وحدة إضافية تكلّف حوالي $36$ ألف دينار، أي $360$ دينارا للوحدة.

\QQ{5}\pa $700$ وحدة توافق $x=7$: \ $y=36\times7+82=252+82=334$. التكلفة \bxt{$334$ ألف دينار تقريبا}
\pb $36x+82\le400$ \ يكافئ \ $x\le\dfrac{318}{36}\approx8{,}83$، أي $883$ وحدة على الأكثر: \bx{883}
\end{hal}

% ------------------------------------------------------------------
\begin{tamrin}{الإشهار ورقم الأعمال -- ماير بقيم غير منتظمة}
يعطي الجدول التالي مصاريف الإشهار $x_i$ (بعشرات الآلاف من الدنانير) ورقم الأعمال $y_i$ (بملايين الدنانير) لمؤسسة خلال ستة أشهر:
\begin{center}\begin{LTR}\renewcommand{\arraystretch}{1.35}\small
\begin{tabular}{|c|c|c|c|c|c|c|}\hline
\rowcolor{cSec!12}\ar{الإشهار} $x_i$ & $2$ & $3$ & $5$ & $6$ & $8$ & $12$\\\hline
\ar{رقم الأعمال} $y_i$ & $2,9$ & $3,8$ & $4,6$ & $5,4$ & $6,1$ & $8,6$\\\hline
\end{tabular}\end{LTR}\end{center}
\QQ{1} مثّل سحابة النقط $M_i(x_i;y_i)$ في معلم متعامد، واحسب إحداثيي $G$.
\QQ{2} $G_1$ و $G_2$ النقطتان المتوسطتان للنقط الثلاث الأولى وللنقط الثلاث الأخيرة.
\begin{sq}
\qa احسب إحداثيي $G_1$ و $G_2$.
\qb بيّن أن معادلة مستقيم ماير $(G_1G_2)$ هي $y=0{,}55x+\dfrac{29}{15}$، وتحقق أنه يشمل $G$.
\end{sq}
\QQ{3}
\begin{sq}
\qa قدّر رقم الأعمال إذا خصصت المؤسسة $150\,000$ دينار للإشهار.
\qb ما مبلغ الإشهار اللازم لبلوغ رقم أعمال قدره $10$ ملايين دينار؟
\end{sq}
\end{tamrin}
\begin{hal}
\QQ{1} السحابة ومستقيم ماير:
\begin{SC}{xmin=0,xmax=13.5,ymin=2,ymax=9.5,xtick={0,2,...,12},ytick={2,...,9},xlabel={$x_i$},ylabel={$y_i$}}
\addplot[ln,domain=0:13.3]{0.55*x+29/15};
\addplot[pt] coordinates {(2,2.9)(3,3.8)(5,4.6)(6,5.4)(8,6.1)(12,8.6)};
\addplot[gq] coordinates {(3.3333,3.7667)(8.6667,6.7)};
\addplot[gp] coordinates {(6,5.2333)};
\node[below right] at (axis cs:6,5.2333){$G$};
\node[above left] at (axis cs:3.3333,3.7667){$G_1$};
\node[below right] at (axis cs:8.6667,6.7){$G_2$};
\node[cThm,above left] at (axis cs:13,9.1){$(G_1G_2)$};
\end{SC}
$\bar x=\dfrac{36}{6}=6$ \ و \ $\bar y=\dfrac{31{,}4}{6}=\dfrac{157}{30}\approx5{,}2333$ \ إذن \bx{G\left(6\,;\,\tfrac{157}{30}\right)}

\QQ{2}\pa \[G_1\left(\frac{2+3+5}{3}\,;\,\frac{2{,}9+3{,}8+4{,}6}{3}\right)=G_1\left(\frac{10}{3}\,;\,\frac{113}{30}\right)\]\[G_2\left(\frac{6+8+12}{3}\,;\,\frac{5{,}4+6{,}1+8{,}6}{3}\right)=G_2\left(\frac{26}{3}\,;\,6{,}7\right)\]
\pb \[a=\frac{6{,}7-\frac{113}{30}}{\frac{26}{3}-\frac{10}{3}}=\frac{\frac{88}{30}}{\frac{16}{3}}=\frac{88}{160}=0{,}55\qquad b=\frac{113}{30}-0{,}55\times\frac{10}{3}=\frac{113-55}{30}=\frac{29}{15}\approx1{,}93\]
إذن \bx{(G_1G_2):\ y=0{,}55x+\tfrac{29}{15}}. \ التحقق: $0{,}55\times6+\dfrac{29}{15}=\dfrac{99+58}{30}=\dfrac{157}{30}=\bar y$ ومنه $G\in(G_1G_2)$ \ok

\QQ{3}\pa $150\,000$ دينار توافق $x=15$: \ $y=0{,}55\times15+\dfrac{29}{15}\approx8{,}25+1{,}93=10{,}18$. رقم الأعمال المقدَّر: \bxt{حوالي $10{,}18$ مليون دينار}
\pb $0{,}55x+\dfrac{29}{15}=10$ \ أي \ $x=\dfrac{10-\frac{29}{15}}{0{,}55}=\dfrac{121/15}{11/20}=\dfrac{44}{3}\approx14{,}67$، أي مبلغ إشهار حوالي \bxt{$146\,700$ دينار}
\end{hal}

% ------------------------------------------------------------------
\begin{tamrin}{مشتركو الإنترنت -- تعديل بياني يمر بـ $G$}
يعطي الجدول عدد المشتركين $y_i$ (بالآلاف) في خدمة الإنترنت عبر الألياف البصرية في بلدية بين 2018 و 2024:
\begin{center}\begin{LTR}\renewcommand{\arraystretch}{1.35}\small
\begin{tabular}{|c|c|c|c|c|c|c|c|}\hline
\rowcolor{cSec!12}\ar{السنة} & $2018$ & $2019$ & $2020$ & $2021$ & $2022$ & $2023$ & $2024$\\\hline
\ar{المشتركون} $y_i$ & $3,2$ & $3,9$ & $4,5$ & $5,4$ & $6$ & $6,8$ & $7,5$\\\hline
\end{tabular}\end{LTR}\end{center}
\QQ{1} نضع $x_i=a_i-2018$ حيث $a_i$ السنة. أعط قيم $x_i$ ثم مثّل سحابة النقط $M_i(x_i;y_i)$.
\QQ{2} احسب إحداثيي النقطة المتوسطة $G$ (دوّر إلى $10^{-4}$).
\QQ{3} نعدّل السحابة بيانيا بالمستقيم $(\Delta)$ الذي يشمل $G$ والنقطة $M_1(0;3{,}2)$.
\begin{sq}
\qa بيّن أن $(\Delta)$ يترك نفس عدد النقط من كل جهة.
\qb عيّن معادلة $(\Delta)$ (دوّر المعامل الموجه إلى $10^{-2}$).
\end{sq}
\QQ{4}
\begin{sq}
\qa قدّر عدد المشتركين سنة 2027.
\qb ابتداء من أي سنة يتجاوز عدد المشتركين $10$ آلاف؟
\end{sq}
\end{tamrin}
\begin{hal}
\QQ{1} السنوات 2018، 2019، $\ldots$، 2024 توافق $x_i=0;1;2;3;4;5;6$:
\begin{SC}{xmin=0,xmax=7.4,ymin=0,ymax=9,xtick={0,...,7},ytick={0,...,9},xlabel={$x_i$},ylabel={$y_i$}}
\addplot[ln,domain=0:7.2]{0.71*x+3.2};
\addplot[pt] coordinates {(0,3.2)(1,3.9)(2,4.5)(3,5.4)(4,6.0)(5,6.8)(6,7.5)};
\addplot[gp] coordinates {(3,5.3286)};
\node[below right] at (axis cs:3,5.3286){$G$};
\node[cThm,below right] at (axis cs:5.6,7.0){$(\Delta)$};
\end{SC}
\QQ{2} $\bar x=\dfrac{21}{7}=3$ ، \ $\bar y=\dfrac{37{,}3}{7}\approx5{,}3286$ \ إذن \bx{G(3\,;\,5{,}3286)}

\QQ{3}\pa على الشكل: $M_1$ على $(\Delta)$؛ النقط ذات الفواصل $3$ ، $5$ ، $6$ فوقه؛ والنقط ذات الفواصل $1$ ، $2$ ، $4$ تحته: \bxt{ثلاث نقط من كل جهة}.

\pb $a=\dfrac{5{,}3286-3{,}2}{3-0}=\dfrac{2{,}1286}{3}\approx0{,}7095\approx0{,}71$ و $(\Delta)$ يقطع محور التراتيب في $3{,}2$: \ \bx{(\Delta):\ y=0{,}71x+3{,}2}

\QQ{4}\pa سنة 2027 توافق $x=9$: \ $y=0{,}71\times9+3{,}2=9{,}59$، أي حوالي \bxt{$9\,590$ مشتركا}
\pb $0{,}71x+3{,}2>10$ \ أي \ $x>\dfrac{6{,}8}{0{,}71}\approx9{,}58$، ومنه $x=10$، أي ابتداء من \bxt{سنة 2028}
\end{hal}

% ------------------------------------------------------------------
\begin{tamrin}{العرض والطلب -- سعر التوازن}
في سوق الجملة، سجّل عدد من المنتجين الكمية المعروضة $y_i$ (بالأطنان) من البطاطا بدلالة سعر الكيلوغرام $x_i$ (بالدينار):
\begin{center}\begin{LTR}\renewcommand{\arraystretch}{1.35}\small
\begin{tabular}{|c|c|c|c|c|c|c|}\hline
\rowcolor{cSec!12}\ar{السعر} $x_i$ & $40$ & $45$ & $50$ & $55$ & $60$ & $65$\\\hline
\ar{الكمية المعروضة} $y_i$ & $120$ & $131$ & $139$ & $152$ & $160$ & $171$\\\hline
\end{tabular}\end{LTR}\end{center}
\QQ{1} مثّل سحابة النقط $M_i(x_i;y_i)$ (المبدأ $(35;110)$)، واحسب إحداثيي $G$.
\QQ{2} باستعمال طريقة ماير (النقط الثلاث الأولى ثم الثلاث الأخيرة)، بيّن أن معادلة مستقيم العرض $(D_1)$ هي $y=2{,}07x+37$ (المعامل الموجه مدوّر إلى $10^{-2}$).
\QQ{3} دراسة أخرى بيّنت أن الكمية المطلوبة تعطى بالمستقيم $(D_2):\ y=-1{,}93x+277$.
\begin{sq}
\qa عيّن سعر التوازن (السعر الذي تتساوى عنده الكميتان المعروضة والمطلوبة) والكمية الموافقة.
\qb إذا حُدّد السعر بـ $65$ DA، احسب الفائض في العرض.
\end{sq}
\end{tamrin}
\begin{hal}
\QQ{1} السحابة، مستقيم العرض $(D_1)$ بالأحمر ومستقيم الطلب $(D_2)$ بالأخضر المتقطع، ونقطة التوازن $E$:
\begin{SC}{xmin=35,xmax=70,ymin=110,ymax=185,xtick={35,40,...,70},ytick={110,120,...,180},xlabel={$x_i$},ylabel={$y_i$}}
\addplot[ln,domain=35:70]{2.07*x+37};
\addplot[lm,domain=35:70]{-1.93*x+277};
\addplot[pt] coordinates {(40,120)(45,131)(50,139)(55,152)(60,160)(65,171)};
\addplot[gp] coordinates {(52.5,145.5)};
\node[above left] at (axis cs:52.5,145.5){$G$};
\addplot[gq] coordinates {(60,161.2)};
\node[above right] at (axis cs:60,161.4){$E$};
\node[cThm,left] at (axis cs:64,176){$(D_1)$};
\node[cRem,below left] at (axis cs:69,145){$(D_2)$};
\end{SC}
\[\bar x=\frac{315}{6}=52{,}5\qquad \bar y=\frac{873}{6}=145{,}5\qquad\bx{G(52{,}5\,;\,145{,}5)}\]
\QQ{2} $G_1\left(\dfrac{40+45+50}{3}\,;\,\dfrac{120+131+139}{3}\right)=G_1(45\,;\,130)$ \ و \ $G_2\left(\dfrac{55+60+65}{3}\,;\,\dfrac{152+160+171}{3}\right)=G_2(60\,;\,161)$.
\[a=\frac{161-130}{60-45}=\frac{31}{15}\approx2{,}0667\approx2{,}07\qquad b=130-\frac{31}{15}\times45=130-93=37\]
إذن \bx{(D_1):\ y=2{,}07x+37}. \ (بالقيمة $a=\frac{31}{15}$: $\frac{31}{15}\times52{,}5+37=108{,}5+37=145{,}5=\bar y$، فالمستقيم يشمل $G$ \ok)

\QQ{3}\pa عند التوازن: $2{,}07x+37=-1{,}93x+277$ \ أي \ $4x=240$ \ أي \ \bx{x=60\ \text{DA}}.

والكمية: $y=2{,}07\times60+37=124{,}2+37=161{,}2$ \ (وفي $(D_2)$: $-115{,}8+277=161{,}2$ \ok)، أي \bxt{$161{,}2$ طنا}

\pb لما $x=65$: العرض $2{,}07\times65+37=171{,}55$ والطلب $-1{,}93\times65+277=151{,}55$. الفائض في العرض: $171{,}55-151{,}55=$ \bxt{$20$ طنا}
\end{hal}

\section{التعديل بتغيير المتغير $z=\ln y$}
% ------------------------------------------------------------------
\begin{tamrin}{طلبيات متجر إلكتروني -- نمو أسي}
يعطي الجدول عدد الطلبيات الشهرية $y_i$ لمتجر إلكتروني خلال الأشهر الستة الأولى من نشاطه ($x_i$ رتبة الشهر):
\begin{center}\begin{LTR}\renewcommand{\arraystretch}{1.35}\small
\begin{tabular}{|c|c|c|c|c|c|c|}\hline
\rowcolor{cSec!12}\ar{الشهر} $x_i$ & $1$ & $2$ & $3$ & $4$ & $5$ & $6$\\\hline
\ar{الطلبيات} $y_i$ & $120$ & $165$ & $230$ & $310$ & $430$ & $590$\\\hline
\end{tabular}\end{LTR}\end{center}
\QQ{1} مثّل سحابة النقط $M_i(x_i;y_i)$. هل شكل السحابة يوحي بتعديل خطي؟
\QQ{2} نضع $z=\ln y$. انقل ثم أتمم الجدول بقيم $z_i$ مدوّرة إلى $10^{-2}$.
\QQ{3}
\begin{sq}
\qa مثّل سحابة النقط $N_i(x_i;z_i)$.
\qb باستعمال طريقة ماير (النقط الثلاث الأولى ثم الثلاث الأخيرة)، عيّن معادلة مستقيم تعديل $z$ بدلالة $x$ (دوّر $a$ و $b$ إلى $10^{-3}$).
\end{sq}
\QQ{4} استنتج أن $y\approx k\,q^{x}$ حيث يطلب تعيين $k$ (إلى $10^{-2}$) و $q$ (إلى $10^{-3}$)، ثم فسّر $q$.
\QQ{5}
\begin{sq}
\qa قدّر عدد الطلبيات في الشهر التاسع.
\qb ابتداء من أي شهر يتجاوز عدد الطلبيات $2000$ طلبية؟
\end{sq}
\end{tamrin}
\begin{hal}
\QQ{1} النقط (بالأزرق) والمنحنى الأسي الذي سنحصل عليه في السؤال 4 (بالأحمر):
\begin{SC}{xmin=0,xmax=7.4,ymin=0,ymax=700,xtick={0,...,7},ytick={0,100,...,700},xlabel={$x_i$},ylabel={$y_i$}}
\addplot[thick,cThm,domain=0:7,samples=60]{88.41*1.372^x};
\addplot[pt] coordinates {(1,120)(2,165)(3,230)(4,310)(5,430)(6,590)};
\end{SC}
النقط \textbf{ليست على استقامة}: الزيادة من شهر لآخر تكبر ($45$ ، $65$ ، $80$ ، $120$ ، $160$)، والسحابة منحنية (تقعر نحو الأعلى). إذن التعديل الخطي \textbf{غير ملائم}، ويوحي الشكل بنمو أسي.

\QQ{2} $z_1=\ln120\approx4{,}79$ ، \ $z_2=\ln165\approx5{,}11$ ، $\ldots$
\begin{center}\begin{LTR}\renewcommand{\arraystretch}{1.35}\small
\begin{tabular}{|c|c|c|c|c|c|c|}\hline
\rowcolor{cSec!12}$x_i$ & $1$ & $2$ & $3$ & $4$ & $5$ & $6$\\\hline
$y_i$ & $120$ & $165$ & $230$ & $310$ & $430$ & $590$\\\hline
$z_i=\ln y_i$ & $4,79$ & $5,11$ & $5,44$ & $5,74$ & $6,06$ & $6,38$\\\hline
\end{tabular}\end{LTR}\end{center}
\QQ{3}\pa السحابة $N_i(x_i;z_i)$ ومستقيم ماير:
\begin{SC}{height=5.6cm,xmin=0,xmax=7.4,ymin=4.4,ymax=6.8,xtick={0,...,7},ytick={4.4,4.8,...,6.8},xlabel={$x_i$},ylabel={$z_i$}}
\addplot[ln,domain=0:7.2]{0.316*x+4.482};
\addplot[gq] coordinates {(2,5.1133)(5,6.06)};
\node[below right] at (axis cs:2,5.1133){$G_1$};
\node[below right] at (axis cs:5,6.06){$G_2$};
\addplot[pt] coordinates {(1,4.79)(2,5.11)(3,5.44)(4,5.74)(5,6.06)(6,6.38)};
\addplot[gp] coordinates {(3.5,5.5867)};
\node[above left] at (axis cs:3.5,5.5867){$G'$};
\end{SC}
النقط $N_i$ تكاد تكون على استقامة واحدة.

\pb المجموعة الأولى ($x_i=1;2;3$) والثانية ($x_i=4;5;6$):
\[G_1\left(2\,;\,\frac{4{,}79+5{,}11+5{,}44}{3}\right)=G_1\left(2\,;\,\frac{15{,}34}{3}\right)\approx G_1(2\,;\,5{,}1133)\]\[G_2\left(5\,;\,\frac{5{,}74+6{,}06+6{,}38}{3}\right)=G_2(5\,;\,6{,}06)\]
\[a=\frac{6{,}06-5{,}1133}{5-2}\approx0{,}3156\approx0{,}316\qquad b=5{,}1133-0{,}3156\times2\approx4{,}482\]
إذن \bx{z=0{,}316x+4{,}482}. \ (تحقق: $\bar z=\dfrac{33{,}52}{6}\approx5{,}5867$ و $0{,}3156\times3{,}5+4{,}4822\approx5{,}5867$، فالمستقيم يشمل $G'$ \ok)

\QQ{4} $z=\ln y$ يعني $y=e^{z}$، ومنه:
\[y=e^{0{,}316x+4{,}482}=e^{4{,}482}\times\left(e^{0{,}316}\right)^{x}\]
مع $e^{4{,}482}\approx88{,}41$ و $e^{0{,}316}\approx1{,}372$. إذن \bx{y\approx88{,}41\times(1{,}372)^{x}}. \ \textbf{التفسير:} $q\approx1{,}372=1+0{,}372$: يزداد عدد الطلبيات كل شهر بنسبة \bx{37{,}2\,\%} تقريبا.

\QQ{5}\pa $x=9$: \ $z=0{,}316\times9+4{,}482=7{,}326$ ومنه $y=e^{7{,}326}\approx1519$. أي حوالي \bxt{$1519$ طلبية}
\pb $y>2000$ \ يكافئ \ $z>\ln2000$ (لأن $\ln$ متزايدة تماما) \ أي \ $0{,}316x+4{,}482>7{,}601$ \ أي \ $x>\dfrac{3{,}119}{0{,}316}\approx9{,}87$.

إذن ابتداء من \bxt{الشهر العاشر}
\end{hal}

% ------------------------------------------------------------------
\begin{tamrin}{سكان مدينة جديدة -- ماير بعد تغيير المتغير}
يعطي الجدول عدد سكان $y_i$ (بالآلاف) لمدينة جديدة بين 2017 و 2024:
\begin{center}\begin{LTR}\renewcommand{\arraystretch}{1.35}\small
\begin{tabular}{|c|c|c|c|c|c|c|c|c|}\hline
\rowcolor{cSec!12}\ar{السنة} & $2017$ & $2018$ & $2019$ & $2020$ & $2021$ & $2022$ & $2023$ & $2024$\\\hline
\ar{الرتبة} $x_i$ & $1$ & $2$ & $3$ & $4$ & $5$ & $6$ & $7$ & $8$\\\hline
\ar{السكان} $y_i$ & $50$ & $53,1$ & $56$ & $59,4$ & $63$ & $66,6$ & $70,6$ & $74,9$\\\hline
\end{tabular}\end{LTR}\end{center}
\QQ{1} نضع $z=\ln y$. أعط قيم $z_i$ مدوّرة إلى $10^{-2}$.
\QQ{2} مثّل سحابة النقط $N_i(x_i;z_i)$. هل يبرّر شكلها تعديلا خطيا؟
\QQ{3} باستعمال طريقة ماير (النقط الأربع الأولى ثم الأربع الأخيرة)، بيّن أن معادلة المستقيم المعدِّل هي \mbox{$z=0{,}058x+3{,}852$} (معاملات مدوّرة إلى $10^{-3}$).
\QQ{4} استنتج عبارة $y$ بدلالة $x$ على الشكل $y=k\,q^{x}$ ، ثم استنتج نسبة النمو السنوي للسكان.
\QQ{5}
\begin{sq}
\qa قدّر عدد سكان المدينة سنة 2028.
\qb في أي سنة يتجاوز عدد السكان $100$ ألف نسمة؟
\end{sq}
\end{tamrin}
\begin{hal}
\QQ{1} $z_i=\ln y_i$ مدوّرة إلى $10^{-2}$:
\begin{center}\begin{LTR}\renewcommand{\arraystretch}{1.35}\small
\begin{tabular}{|c|c|c|c|c|c|c|c|c|}\hline
\rowcolor{cSec!12}$x_i$ & $1$ & $2$ & $3$ & $4$ & $5$ & $6$ & $7$ & $8$\\\hline
$y_i$ & $50$ & $53,1$ & $56$ & $59,4$ & $63$ & $66,6$ & $70,6$ & $74,9$\\\hline
$z_i=\ln y_i$ & $3,91$ & $3,97$ & $4,03$ & $4,08$ & $4,14$ & $4,2$ & $4,26$ & $4,32$\\\hline
\end{tabular}\end{LTR}\end{center}
\QQ{2} السحابة $N_i(x_i;z_i)$ مع مستقيم ماير:
\begin{SC}{height=5.8cm,xmin=0,xmax=9.3,ymin=3.85,ymax=4.4,xtick={0,...,9},ytick={3.85,3.95,...,4.35},yticklabel style={/pgf/number format/fixed,/pgf/number format/precision=2,/pgf/number format/use comma},xlabel={$x_i$},ylabel={$z_i$}}
\addplot[ln,domain=0:9.2]{0.058*x+3.852};
\addplot[pt] coordinates {(1,3.91)(2,3.97)(3,4.03)(4,4.08)(5,4.14)(6,4.20)(7,4.26)(8,4.32)};
\addplot[gq] coordinates {(2.5,3.9975)(6.5,4.23)};
\node[below right] at (axis cs:2.5,3.9975){$G_1$};
\node[below right] at (axis cs:6.5,4.23){$G_2$};
\end{SC}
النقط $N_i$ تقريبا على استقامة واحدة ($z$ يزداد بحوالي $0{,}06$ كل سنة)، إذن التعديل الخطي لـ $z$ بدلالة $x$ مبرَّر.

\QQ{3} \textbf{المجموعة الأولى} ($x_i=1;2;3;4$):
\[x_{G_1}=\frac{1+2+3+4}{4}=2{,}5\qquad z_{G_1}=\frac{3{,}91+3{,}97+4{,}03+4{,}08}{4}=\frac{15{,}99}{4}=3{,}9975\]
\textbf{المجموعة الثانية} ($x_i=5;6;7;8$):
\[x_{G_2}=\frac{5+6+7+8}{4}=6{,}5\qquad z_{G_2}=\frac{4{,}14+4{,}20+4{,}26+4{,}32}{4}=\frac{16{,}92}{4}=4{,}23\]
\[a=\frac{4{,}23-3{,}9975}{6{,}5-2{,}5}=\frac{0{,}2325}{4}=0{,}058125\approx0{,}058\qquad b=3{,}9975-0{,}058125\times2{,}5\approx3{,}852\]
إذن \bx{z=0{,}058x+3{,}852}. \ (تحقق: $\bar z=\dfrac{32{,}91}{8}\approx4{,}114$ و $0{,}058125\times4{,}5+3{,}8522\approx4{,}114$ \ok)

\QQ{4} $y=e^{z}=e^{3{,}852}\times\left(e^{0{,}058}\right)^{x}$ \ مع \ $e^{3{,}852}\approx47{,}09$ و $e^{0{,}058}\approx1{,}06$. إذن \bx{y\approx47{,}09\times(1{,}06)^{x}}

$q\approx1{,}06=1+0{,}06$: يزداد عدد السكان كل سنة بنسبة \bx{6\,\%} تقريبا.

\QQ{5}\pa سنة 2028 توافق $x=2028-2016=12$: \ $z=0{,}058\times12+3{,}852=4{,}548$ ومنه $y=e^{4{,}548}\approx94{,}44$، أي حوالي \bxt{$94\,440$ نسمة}
\pb $y>100$ \ يكافئ \ $z>\ln100\approx4{,}605$ \ أي \ $0{,}058x>0{,}753$ \ أي \ $x>12{,}98$ تقريبا، ومنه $x=13$، أي \bxt{سنة 2029}
\end{hal}

% ------------------------------------------------------------------
\begin{tamrin}{اهتلاك سيارة -- نقصان أسي}
يعطي الجدول قيمة سيارة $y_i$ (بملايين الدنانير) بدلالة عمرها $x_i$ (بالسنوات) في سوق السيارات المستعملة:
\begin{center}\begin{LTR}\renewcommand{\arraystretch}{1.35}\small
\begin{tabular}{|c|c|c|c|c|c|c|c|}\hline
\rowcolor{cSec!12}\ar{العمر} $x_i$ & $0$ & $1$ & $2$ & $3$ & $4$ & $5$ & $6$\\\hline
\ar{القيمة} $y_i$ & $2,4$ & $2,04$ & $1,75$ & $1,49$ & $1,27$ & $1,08$ & $0,93$\\\hline
\end{tabular}\end{LTR}\end{center}
\QQ{1} مثّل سحابة النقط $M_i(x_i;y_i)$ وصف شكلها.
\QQ{2} نضع $z=\ln y$. أتمم جدول القيم بـ $z_i$ مدوّرة إلى $10^{-2}$، ثم احسب إحداثيي النقطة المتوسطة $G'$ للسلسلة $(x_i;z_i)$.
\QQ{3} مثّل السحابة $N_i(x_i;z_i)$. أي التعديلين أنسب: $y$ بدلالة $x$ أم $z$ بدلالة $x$؟
\QQ{4} نقبل أن المستقيم $(\Delta):\ z=-0{,}158x+0{,}875$ مستقيم تعديل لـ $z$ بدلالة $x$. تحقق أنه يشمل $G'$ (بتقريب $10^{-2}$).
\QQ{5} استنتج أن $y\approx2{,}40\times(0{,}854)^{x}$، ثم استنتج نسبة الاهتلاك السنوي.
\QQ{6}
\begin{sq}
\qa قدّر قيمة السيارة بعد $8$ سنوات.
\qb ابتداء من أي عمر تصبح قيمة السيارة أقل من نصف قيمتها وهي جديدة ($2{,}40$ مليون دينار)؟
\end{sq}
\end{tamrin}
\begin{hal}
\QQ{1} السحابة والمنحنى الأسي المعدِّل (السؤال 5):
\begin{SC}{height=5.8cm,xmin=0,xmax=8.4,ymin=0,ymax=2.6,xtick={0,...,8},ytick={0,0.5,...,2.5},xlabel={$x_i$},ylabel={$y_i$}}
\addplot[thick,cThm,domain=0:8.2,samples=60]{2.40*0.854^x};
\addplot[pt] coordinates {(0,2.40)(1,2.04)(2,1.75)(3,1.49)(4,1.27)(5,1.08)(6,0.93)};
\end{SC}
القيمة تتناقص، لكن النقصان السنوي يصغر ($0{,}36$ ، $0{,}29$ ، $0{,}26$ ، $0{,}22$ ، $0{,}19$ ، $0{,}15$): السحابة منحنية وليست على استقامة.

\QQ{2} قيم $z_i$:
\begin{center}\begin{LTR}\renewcommand{\arraystretch}{1.35}\small
\begin{tabular}{|c|c|c|c|c|c|c|c|}\hline
\rowcolor{cSec!12}$x_i$ & $0$ & $1$ & $2$ & $3$ & $4$ & $5$ & $6$\\\hline
$y_i$ & $2,4$ & $2,04$ & $1,75$ & $1,49$ & $1,27$ & $1,08$ & $0,93$\\\hline
$z_i=\ln y_i$ & $0,88$ & $0,71$ & $0,56$ & $0,4$ & $0,24$ & $0,08$ & $-0,07$\\\hline
\end{tabular}\end{LTR}\end{center}
\[\bar x=\frac{21}{7}=3\qquad \bar z=\frac{0{,}88+0{,}71+0{,}56+0{,}4+0{,}24+0{,}08-0{,}07}{7}=\frac{2{,}8}{7}=0{,}4\]
إذن \bx{G'(3\,;\,0{,}4)}

\QQ{3} السحابة $N_i(x_i;z_i)$ والمستقيم $(\Delta)$:
\begin{SC}{height=5.6cm,xmin=0,xmax=7.4,ymin=-0.2,ymax=1,xtick={0,...,7},ytick={-0.2,0,...,1},yticklabel style={/pgf/number format/fixed,/pgf/number format/precision=1,/pgf/number format/use comma},xlabel={$x_i$},ylabel={$z_i$}}
\addplot[ln,domain=0:7.2]{-0.158*x+0.875};
\addplot[pt] coordinates {(0,0.88)(1,0.71)(2,0.56)(3,0.4)(4,0.24)(5,0.08)(6,-0.07)};
\addplot[gp] coordinates {(3,0.4)};
\node[above right] at (axis cs:3,0.4){$G'$};
\end{SC}
النقط $N_i$ تكاد تكون على استقامة ($z$ ينقص بحوالي $0{,}16$ كل سنة)، بينما السحابة $M_i$ منحنية: تعديل \textbf{$z$ بدلالة $x$} أنسب.

\QQ{4} $-0{,}158\times3+0{,}875=-0{,}474+0{,}875=0{,}401\approx0{,}4=\bar z$، إذن \bx{G'\in(\Delta)} (بتقريب $10^{-2}$) \ok

\QQ{5} $y=e^{z}=e^{0{,}875}\times\left(e^{-0{,}158}\right)^{x}$ \ مع \ $e^{0{,}875}\approx2{,}40$ و $e^{-0{,}158}\approx0{,}854$. إذن \bx{y\approx2{,}40\times(0{,}854)^{x}}

$q\approx0{,}854=1-0{,}146$: تفقد السيارة كل سنة حوالي \bx{14{,}6\,\%} من قيمتها.

\QQ{6}\pa $x=8$: \ $z=-0{,}158\times8+0{,}875=-0{,}389$ ومنه $y=e^{-0{,}389}\approx0{,}678$، أي حوالي \bxt{$678\,000$ دينار}
\pb $y<1{,}20$ \ يكافئ \ $z<\ln1{,}2\approx0{,}182$ \ أي \ $-0{,}158x<-0{,}693$ \ أي \ $x>\dfrac{0{,}693}{0{,}158}\approx4{,}39$.

إذن ابتداء من \bxt{عمر $5$ سنوات} تصبح قيمتها أقل من نصف قيمتها الأصلية (يتوافق مع الجدول: $1{,}27$ عند 4 سنوات و $1{,}08$ عند 5 سنوات).
\end{hal}

\section{تمرين تقييمي}
% ------------------------------------------------------------------
\begin{tamrin}{اختيار من متعدد وصحيح أم خاطئ}
نعتبر السلسلة الإحصائية ذات المتغيرين:
\begin{center}\begin{LTR}\renewcommand{\arraystretch}{1.35}\small
\begin{tabular}{|c|c|c|c|c|c|c|}\hline
\rowcolor{cSec!12}$x_i$ & $1$ & $2$ & $3$ & $4$ & $5$ & $6$\\\hline
$y_i$ & $3$ & $5$ & $4$ & $7$ & $6$ & $8$\\\hline
\end{tabular}\end{LTR}\end{center}
\textbf{الجزء الأول:} اختر الجواب الصحيح الوحيد في كل حالة.
\QQ{1} النقطة المتوسطة $G$ هي: \ABC{$G(3{,}5;5{,}5)$}{$G(5{,}5;3{,}5)$}{$G(3;5{,}5)$}
\QQ{2} النقطة المتوسطة $G_1$ للنقط الثلاث الأولى هي: \ABC{$G_1(2;4)$}{$G_1(2;5)$}{$G_1(4;2)$}
\QQ{3} مستقيم ماير $(G_1G_2)$ (النقط الثلاث الأولى ثم الثلاث الأخيرة) هو: \ABC{$y=x+2$}{$y=2x+1$}{$y=x+3{,}5$}
\QQ{4} حسب مستقيم ماير، تبلغ القيمة المقدَّرة لـ $y$ العدد $20$ لما: \ABC{$x=18$}{$x=22$}{$x=10$}

\textbf{الجزء الثاني:} أجب بصحيح أو خاطئ مع التعليل.
\QQ{5} مستقيم ماير $(G_1G_2)$ يشمل النقطة المتوسطة $G$.
\QQ{6} المستقيم $(\Delta):\ y=1{,}2x+1{,}5$ يشمل النقطة $G$.
\QQ{7} إذا أضفنا $10$ إلى كل القيم $y_i$ فإن $G$ تصبح $G(3{,}5;15{,}5)$ ومعامل توجيه مستقيم ماير لا يتغير.
\QQ{8} القيمة المقدَّرة لـ $y$ من أجل $x=10$ بمستقيم ماير هي $13$.
\end{tamrin}
\begin{hal}
$\sum x_i=21$ ، $\sum y_i=3+5+4+7+6+8=33$.

\QQ{1} $\bar x=\dfrac{21}{6}=3{,}5$ و $\bar y=\dfrac{33}{6}=5{,}5$: الجواب \bxt{(أ) $G(3{,}5;5{,}5)$}
\QQ{2} $x_{G_1}=\dfrac{1+2+3}{3}=2$ و $y_{G_1}=\dfrac{3+5+4}{3}=4$: الجواب \bxt{(أ) $G_1(2;4)$}
\QQ{3} $G_2\left(\dfrac{4+5+6}{3};\dfrac{7+6+8}{3}\right)=G_2(5;7)$، $a=\dfrac{7-4}{5-2}=1$ و $b=4-1\times2=2$: الجواب \bxt{(أ) $y=x+2$}
\QQ{4} $x+2=20$ يعني $x=18$: الجواب \bxt{(أ)}
\QQ{5} \bxt{صحيح}: $3{,}5+2=5{,}5=\bar y$ (المجموعتان لهما نفس عدد النقط).
\QQ{6} \bxt{خاطئ}: $1{,}2\times3{,}5+1{,}5=4{,}2+1{,}5=5{,}7\neq5{,}5$.
\QQ{7} \bxt{صحيح}: $\bar y$ و $y_{G_1}$ و $y_{G_2}$ تزداد كلها بـ $10$، فـ $G(3{,}5;15{,}5)$ والفرق $y_{G_2}-y_{G_1}$ لا يتغير، أي $a=1$ دائما.
\QQ{8} \bxt{خاطئ}: $10+2=12\neq13$.
\end{hal}
\end{document}
